\chapter{YAPAY ZEKÂ VE METODOLOJİ}

Aura Finance, kullanıcıların finansal kararlarını desteklemek için iki temel yapay zekâ yaklaşımı kullanmaktadır: Zaman serisi tahmini (LSTM) ve bütçe optimizasyonu (XGBoost).

\section{Harcama Tahmini: LSTM Modeli}

Kullanıcıların gelecek ay ne kadar harcama yapacağını tahmin etmek için Uzun Kısa Süreli Bellek (Long Short-Term Memory - LSTM) ağları kullanılmıştır. Bu model, Node.js üzerinde \texttt{TensorFlow.js} kütüphanesi ile gerçek zamanlı olarak çalıştırılmaktadır.

\subsection{Model Mimarisi}
Model, ardışık (sequential) bir yapıdadır ve şu katmanlardan oluşur: LSTM katmanı (32 birim, \%10 dropout), Yoğun katman (16 birim, ReLU) ve Çıkış katmanı (1 birim, sigmoid).

\subsection{Eğitim Süreci}
Model, kullanıcının geçmiş harcama verilerini pencereleme (windowing) yöntemiyle işler. Veriler eğitim öncesinde [0, 1] aralığına normalleştirilir. Adam optimizasyon algoritması kullanılarak eğitilir.

\section{Bütçe Optimizasyonu: XGBoost Regresyonu}

Kullanıcıların gelirlerine göre ideal bütçe dağılımını belirlemek için Python tabanlı bir makine öğrenmesi modeli geliştirilmiştir. Eğitim için \texttt{Indian\_Finance\_Dataset.csv} kullanılmıştır. Model olarak \texttt{MultiOutputRegressor} mimarisiyle XGBoost algoritması tercih edilmiştir.

\section{Hibrit Dürtme (Nudge) Algoritması}

Sistemin en özgün yanlarından biri, statik bütçe önerileri yerine kullanıcının mevcut alışkanlıklarını da hesaba katan hibrit bir mantık yürütmesidir. İdeal bütçe ($B_{ideal}$) ve kullanıcının gerçek harcama ortalaması ($B_{gercek}$) arasındaki denge şu formülle hesaplanır:

\begin{equation}
   B_{oneri} = (B_{ideal} \times 0.70) + (B_{gercek} \times 0.30)
\end{equation}

Bu yaklaşım, kullanıcıyı ideal finansal hedeflere yönlendirirken, mevcut yaşam standartlarını da göz ardı etmeyerek gerçekçi bütçe hedefleri sunar.
